\section{Parte 5: pronóstico y fuga}
\label{sec:parte5}

Cada modelo está implementado en \texttt{parte5.py} como una función que
recibe la serie, el origen (un domingo) y los días abiertos de la semana
siguiente, y que solo puede mirar la serie hasta el origen. El origen móvil
sigue exactamente las reglas del enunciado: los objetivos son los días
abiertos de lunes a domingo después del origen; en (b) no se evalúa nada
después del 30 de septiembre y en (c) el último día es el 8 de diciembre; para
S1 y S2 el horizonte es la posición del objetivo entre los días abiertos
posteriores al origen; y el MAE y el RMSE juntan todos los días pronosticados.
En (b) hay 22 orígenes, del 1 de mayo al 25 de septiembre, con 128 días
evaluados; en (c) hay 10, del 2 de octubre al 4 de diciembre, con 58.

\subsection{Inciso 5.1: los modelos}

\begin{itemize}
  \item \textbf{B1}: el promedio de todos los días abiertos hasta el origen.
  \item \textbf{B2} (\emph{naive} estacional): para cada objetivo, la venta del
    mismo día de la semana abierto más reciente hasta el origen.
  \item \textbf{M1}: OLS de $y_t$ contra $t$ (días desde el 1 de diciembre de
    2010) y las seis indicadoras de día de la semana.
  \item \textbf{M2}: OLS de $y_t$ contra las seis indicadoras, $y_{t-7}$ y
    $y_{t-14}$. Si el día del rezago está cerrado se usa el mismo día de la
    semana abierto más reciente; el código verifica con un \texttt{assert} el
    ejemplo del enunciado (para el 9 de mayo, $y_{t-7}$ es la venta del 18 de
    abril, \pounds31\,402). Para $y_{t-14}$ se aplica la misma regla, así que
    el 9 de mayo sus dos rezagos son el 18 de abril. Las 12 filas del inicio
    que no tienen rezago no se usan para ajustar.
  \item \textbf{S1} y \textbf{S2}: los modelos del inciso 4.3, reajustados en
    cada origen con todos los días abiertos hasta el origen y llevados a
    libras con $\exp(\cdot)$.
\end{itemize}

\paragraph{¿Por qué M2 puede usar $y_{t-7}$ pero no $y_{t-1}$?} El
pronóstico se hace el domingo para toda la semana siguiente. Para cualquier
objetivo de lunes a domingo, siete días antes cae en el domingo de origen o
antes, así que $y_{t-7}$ ya se observó: el lunes usa el lunes anterior y el
domingo usa el propio domingo de origen. En cambio, $y_{t-1}$ solo se conoce
para el lunes; para el martes, $y_{t-1}$ es el lunes, que el domingo todavía
no ha ocurrido. Una regresión que lo usara como variable necesitaría un dato
del futuro.

\paragraph{¿Cómo usa S1 el rezago 1 de todas formas?} De forma recursiva. Un
modelo autorregresivo pronostica un paso a la vez y, cuando le falta un
valor, usa su propio pronóstico en su lugar: $\hat y_{t+2} = c + \phi\,\hat
y_{t+1} + \dots$. El filtro de Kalman de \texttt{SARIMAX} hace exactamente
eso. Por eso la información del último día observado se diluye como $\phi^h$
conforme crece el horizonte $h$: con $\phi\approx0.13$, para el martes ya casi
no queda nada del domingo. M2 no puede hacer esto porque es una regresión
directa: no pronostica sus propias variables.

\paragraph{Convergencia en cada origen.} Se reajustan S1 y S2 en los 32
orígenes de (b) y (c), con la configuración del inciso 4.3
(\texttt{maxiter=500}), y se revisa \texttt{mle\_retvals['converged']} en
cada uno. S2 convergió en los 32. S1 convergió en 31; en el del 1 de mayo
L-BFGS se detuvo en la iteración 52 con \texttt{warnflag=2}, que significa que
la búsqueda de línea ya no lograba mejorar. No fue por falta de iteraciones:
con 1\,000 o 2\,000 se detiene en el mismo punto, con un gradiente del orden
de $10^{-4}$. El código reintenta con BFGS, que sí converge, a los mismos
parámetros hasta el segundo decimal y con log-verosimilitud $-80.3911$ contra
$-80.3912$. Los pronósticos cambian menos de \pounds12. Es decir, el ajuste
estaba en el óptimo y el optimizador no lo pudo certificar; aun así se usa el
ajuste que sí convergió.

\subsection{Inciso 5.2: tres formas de estimar el error}

La tabla~\ref{tab:p5_errores} reúne los tres esquemas. En (a), cada celda es
la media $\pm$ la desviación estándar sobre las 10 semillas
(\texttt{SEMILLAS\_CV} en \texttt{semillas.py}) de una validación cruzada
aleatoria de 5 \emph{folds} sobre los días abiertos del entrenamiento.

\begin{table}[H]
  \centering
  \small
  \tabla{p5_errores}
  \caption{MAE y RMSE en libras: (a) validación cruzada aleatoria de 5
  \emph{folds} en el entrenamiento, con 10 semillas; (b) origen móvil semanal
  en el entrenamiento, orígenes de mayo a septiembre; (c) origen móvil
  semanal en la prueba.}
  \label{tab:p5_errores}
\end{table}

\paragraph{¿Por qué no tiene sentido (a) para B2, S1 y S2?} Porque los tres
dependen del orden de la serie, y la validación cruzada aleatoria lo destruye.
Al sacar al azar un 20\% de los días, el día a pronosticar queda rodeado de
días de entrenamiento \emph{posteriores} a él. Para B2, ``el mismo día de la
semana más reciente'' de un día de marzo podría ser un día de abril, y eso ya
no es un pronóstico sino una interpolación. S1 y S2 necesitan la secuencia
completa para que $B$ y $B^6$ signifiquen ``ayer'' y ``hace una semana'': con
huecos aleatorios el rezago $k$ deja de ser un rezago fijo, y el filtro, al
estimar un día faltante, usaría los días siguientes. B1, M1 y M2 son
regresiones cuyas filas se pueden repartir entre \emph{folds}, así que
mecánicamente sí se puede calcular (a) para ellas. Pero también están
contaminadas: M1 estima la pendiente con días posteriores al que pronostica, y
el nivel de los días vecinos (que son los que más se parecen) siempre queda
en el entrenamiento.

\paragraph{¿Cuál se acerca más a (c)?} El origen móvil. Para M1 y M2, (b)
da errores más grandes que (a) y, por lo tanto, más cercanos a (c): M2 pasa
de 8\,345 en (a) a 9\,258 en (b), contra 15\,219 en (c), y M1 de 9\,129 a
9\,833, contra 15\,203. Para B1 los dos esquemas coinciden (10\,271 y
10\,253): la media histórica no tiene cómo aprovechar a los vecinos de un
día, así que la validación aleatoria no le regala nada. La validación cruzada aleatoria es
la más optimista porque es la que más información del futuro deja entrar.

\paragraph{¿Alguna se acerca de verdad?} No. El error en la prueba es entre
1.25 y 2.2 veces el de (b), y además el orden de los modelos cambia por
completo. En (b) el mejor es M2 y S2 queda cuarto, por delante solo de B1 y
B2; en (c) S2 es claramente el mejor y M2 queda en el lugar cuatro, empatado
con M1. Ninguna estimación hecha dentro del entrenamiento habría escogido el
modelo ganador. El inciso 5.5 discute por qué.

\subsection{Inciso 5.3: qué modelo gana en la prueba}

\textbf{Gana S2}, con MAE de \pounds12\,545 y RMSE de \pounds15\,901: un 13\%
menos de MAE que el segundo lugar (B2, \pounds14\,409) y un 18\% menos que M1
y M2. La figura~\ref{fig:p5_prueba} muestra los pronósticos de todos los
modelos contra la serie real.

\begin{figure}[H]
  \centering
  \includegraphics[width=\textwidth]{p5_pronosticos_prueba.png}
  \caption{Pronósticos con origen móvil semanal en la prueba (octubre a
  diciembre de 2011) contra la venta real. Cada línea punteada es un domingo
  de origen; cada modelo pronostica los seis días abiertos que le siguen.}
  \label{fig:p5_prueba}
\end{figure}

\begin{table}[H]
  \centering
  \small
  \tabla{p5_sesgo}
  \caption{Error medio en la prueba (real menos pronóstico), en libras. Un
  valor positivo significa que el modelo se queda corto.}
  \label{tab:p5_sesgo}
\end{table}

\paragraph{Por qué gana S2.} La prueba es la temporada navideña: la venta
media de los días evaluados en (c) es \pounds49\,800, contra \pounds29\,500
en (b). La tabla~\ref{tab:p5_sesgo} muestra qué hace cada modelo con ese
cambio de nivel. B1, M1, M2 y S1 se quedan cortos, en promedio, por
\pounds13\,000 a \pounds20\,500 diarios. B1 promedia todo el año. M1 ajusta
una sola recta a un año que fue plano ocho meses y luego subió, y la recta
queda por debajo del final. M2 le da poco peso a los rezagos (el día de la
semana explica casi todo en el entrenamiento). Y S1, como se vio en el inciso
4.5, regresa a una media estimada con enero a septiembre. En la
figura~\ref{fig:p5_prueba} se ve que sus pronósticos tienen la forma semanal
correcta pero quedan sistemáticamente por debajo de la línea negra a partir
de noviembre.

Solo dos modelos casi no tienen sesgo: B2 (\pounds2\,743) y S2
(\pounds1\,187). Los dos parten de lo que pasó la semana anterior y, por lo
tanto, siguen el nivel. La diferencia entre ellos es el ruido: B2 copia una
sola semana, con todos sus picos, así que cuando el martes anterior tuvo un
pedido excepcional lo repite (para el lunes 21 de noviembre pronostica
\pounds110 mil porque ese fue el pico real del lunes 14; la venta fue de
\pounds47 mil). S2 hace lo mismo pero promediando: sus
términos de promedio móvil, $\theta=-0.91$ y $\Theta=-0.67$, convierten la
diferencia en un promedio exponencialmente ponderado de las semanas
recientes. Sigue el nivel como B2, pero sin copiar el ruido.

Esto es coherente con el diagnóstico del inciso 4.4: S2 fue el único modelo
cuyos residuos parecían ruido blanco, porque fue el único que capturó el
nivel que se mueve. S1 dejaba autocorrelación en los residuos, justamente por
ese nivel, y en la prueba eso se convierte en sesgo. En (b), de mayo a
septiembre, el nivel fue estable y modelos que regresan a la media no pagaban
ese costo; la flexibilidad de S2 solo le añadía varianza. En la prueba, el
nivel se movió y esa flexibilidad es la que gana.

\begin{table}[H]
  \centering
  \small
  \tabla{p5_semanas}
  \caption{MAE de cada semana de la prueba, en miles de libras. ``Mejor'' es
  el modelo con menor MAE esa semana.}
  \label{tab:p5_semanas}
\end{table}

La victoria no es de todas las semanas (tabla~\ref{tab:p5_semanas}): S2 es el
mejor en 5 de las 10 y M1 en 3, siempre en semanas tranquilas, sin cambio de
nivel. Pero S2 nunca tiene una semana desastrosa, y por eso gana en el total.
Los errores más grandes de todos los modelos vienen de días con pedidos
excepcionales, como el lunes 14 de noviembre (\pounds110 mil), que ningún
modelo podía anticipar.

\subsection{Inciso 5.4: la trampa}

Se agrega a M2 la media móvil centrada de 7 días de la Parte 3, calculada una
sola vez sobre toda la serie, como una variable más.

\begin{table}[H]
  \centering
  \small
  \tabla{p5_trampa}
  \caption{M2 con y sin la media móvil centrada como predictor. (a):
  promedio sobre 10 semillas; (b): origen móvil en el entrenamiento.}
  \label{tab:p5_trampa}
\end{table}

\paragraph{Se ve mejor.} Mucho mejor: el MAE baja de 8\,345 a 7\,450 en (a) y
de 9\,258 a 8\,377 en (b), alrededor de un 10\% en los dos esquemas
(tabla~\ref{tab:p5_trampa}). El ajuste con todo el entrenamiento le da a la
media móvil un coeficiente de 1.05 y deja los rezagos en prácticamente cero
(0.01 y 0.05): el modelo descubre que la mejor forma de pronosticar la venta
de un día es ``la media móvil de ese día más el efecto del día de la semana''.

\paragraph{Por qué es ilegal.} Porque $MA_t$ contiene el futuro. La media
centrada del día $t$ promedia de $y_{t-3}$ a $y_{t+3}$, así que incluye la
venta del propio día que se quiere pronosticar ($y_t$, con peso $1/7$) y la
de los tres días siguientes. El domingo de origen, para pronosticar el lunes
se necesitarían las ventas de lunes a jueves; para pronosticar el jueves,
hasta las del domingo siguiente. Ningún dato de la semana objetivo existe
cuando se hace el pronóstico.

\paragraph{Por qué el origen móvil no lo detecta.} Porque el origen móvil solo
controla qué \emph{filas} entran al ajuste; no revisa cómo se calcularon las
\emph{variables} de esas filas. Si la media móvil se calcula una vez sobre toda
la serie y luego se parte en orígenes, la columna \texttt{ma7} del lunes
objetivo ya trae adentro las ventas de la semana objetivo, aunque esas filas
no se hayan usado para estimar los coeficientes. La fuga está en la
construcción de la variable, antes de la validación, y ninguna partición de
filas la puede ver. La única defensa es recalcular cada variable en cada
origen usando solo datos hasta el origen, y con la media centrada eso es
imposible para la semana objetivo.

\paragraph{Por qué (c) no se puede calcular sin hacer trampa.} En la prueba,
para el origen del 2 de octubre la media móvil de los objetivos necesita las
ventas del 3 al 11 de octubre, que son precisamente los valores que se están
evaluando. Calcular (c) con esa variable es darle al modelo parte de la
respuesta del examen. La versión honesta requeriría pronosticar primero la
media móvil, es decir, pronosticar las ventas de la semana, que es justo el
problema que se quería resolver. Por eso esta variable no tiene versión
legítima para pronosticar, y por eso su mejora en (a) y (b) es una ilusión.

\subsection{Inciso 5.5: por qué el entrenamiento se queda corto}

Todas las estimaciones hechas dentro del entrenamiento subestiman el error de
la prueba. Hay dos causas, y conviene separarlas.

\paragraph{1. Cómo se partieron los datos.} La validación cruzada aleatoria
deja que el modelo aprenda del futuro y que el día pronosticado esté rodeado
de días vecinos en el entrenamiento, que comparten su nivel. Es lo que separa
(a) de (b): de 8\,345 a 9\,258 en M2, de 9\,129 a 9\,833 en M1. \emph{Esta
causa sí se arregla con una mejor validación}: el origen móvil reproduce las
condiciones reales del pronóstico (solo pasado, horizonte de una semana) y la
elimina.

\paragraph{2. La serie cambió.} La prueba es la temporada navideña, con una
venta media de \pounds49\,800 contra \pounds29\,500 de los días evaluados en
(b), y los errores crecen con el nivel. Si el MAE se divide entre la venta
media de los días evaluados, los modelos tienen errores relativos parecidos
en los dos periodos: entre 31\% y 38\% en (b), y entre 25\% y 45\% en (c). M2,
por ejemplo, tiene un 31\% en ambos. Lo que más cambia es para quién: B1 y
S1, que regresan a la media, empeoran en términos relativos; S2 y B2, que
siguen el nivel, mejoran. \emph{Esta causa no se arregla con ninguna
validación}: no existe forma de que una validación dentro de enero a
septiembre mida cómo se comporta un modelo en un régimen que el entrenamiento
nunca contuvo. Lo único que ayuda es tener datos de ese régimen, o un modelo
que se adapte a cambios de nivel (como S2), y aun así sin poder medir de
antemano cuánto.

\paragraph{¿Se puede estimar la estacionalidad anual con un año de datos?}
No. Para estimar un patrón que se repite cada año hay que verlo repetirse,
es decir, al menos dos años completos; FPP2 lo plantea igual para cualquier
estacionalidad. Con un año, cada día del año aparece una sola vez, y el
``efecto noviembre'' es indistinguible de una tendencia, de un cambio en la
clientela o de un evento único: los tres producen la misma subida desde
septiembre. Además, el único otro diciembre de los datos (el de 2010) está
partido por el cierre de Navidad. Es justamente lo que le faltó a todos los
modelos: ninguno sabía que el cuarto trimestre vende el doble, porque en el
entrenamiento no había un cuarto trimestre del cual aprenderlo.
